\subsection*{Parte 3: Examen rápido, ¿caja o fracaso?}
Python acepta \texttt{vidas}, \texttt{jugador\_1} y \texttt{mejorTiempo}. \texttt{mejor marca} tiene un espacio, \texttt{1er\_lugar} empieza con un dígito, e \texttt{import} es una palabra que Python ya usa.

\subsection*{Parte 3, Misión: Contador de puntos}
\py|puntos += 10| hace exactamente lo mismo que \py|puntos = puntos + 10|, así que nada cambia. Para contar hacia atrás, empieza con \py|puntos = 50| y usa \py|puntos -= 10|: la última línea dice \texttt{Puntos finales: 0}. Con \py|vidas -= 1| dentro del bucle, las vidas van 2, 1, 0, luego $-1$ y $-2$. Python no se detiene en 0 a menos que tú se lo digas (aprenderás cómo con \py|if| en la Parte 5).

\textbf{Superreto:}
\begin{codenonum}
import time
puntos = 0
premio = 10
print("Puntos:", puntos)
for i in range(5):
    time.sleep(1)
    puntos = puntos + premio
    premio = premio + 10
    print("Puntos:", puntos)
print("¡Fin del juego! Puntos finales:", puntos)
\end{codenonum}
Los puntos van 10, 30, 60, 100, 150. Los puntos finales son \texttt{150}.

\subsection*{Parte 3, Misión: Espiral que crece}
Con \texttt{+ 8} la última línea mide 317 pasos, y la espiral apenas cabe en el escenario. Para el lápiz más grueso, pon \py|grueso = 1| junto a \py|medida = 5|, y agrega \py|pensize(grueso)| y \py|grueso = grueso + 0.2| dentro del bucle.

\textbf{Superreto} (después de las tres líneas de inicio):
\begin{codenonum}
speed(13)
pencolor("purple")
penup()
goto(-100, 100)
pendown()
medida = 200
for i in range(40):
    forward(medida)
    right(90)
    medida = medida - 5
\end{codenonum}

\subsection*{Parte 3, Misión: Arcoíris con una lista de colores}
\py|colores.append("pink")| añade el rosa al final de la lista, así que el bucle da 7 vueltas y el rosa es la franja de hasta abajo. Las franjas roja, blanca y azul, de arriba abajo, forman la bandera de los Países Bajos. Prueba \py|["white", "blue", "red"]| para Rusia.

\textbf{Superreto:} las líneas 8 a 16 quedan así
\begin{codenonum}
y = 60
largo = 300
for color in colores:
    pencolor(color)
    penup()
    goto(-150, y)
    pendown()
    forward(largo)
    y = y - 15
    largo = largo - 30
hideturtle()
\end{codenonum}

\subsection*{Parte 3, Misión: Inventor de historias al azar}
Para los lugares, agrega \py|lugares = ["en la tina", "en el camión"]| antes del bucle, \py|lugar = random.choice(lugares)| dentro de él, y cambia el print por \py|print(heroe, accion, cosa, lugar)|. Después del bucle, \py|len(cosas)| es \texttt{4}, y \texttt{"mi tarea"} es el último elemento, porque \py|append| siempre añade al final:
\begin{outputbox}
4\\
{[}'un plátano', 'la luna', 'un calcetín', 'mi tarea']
\end{outputbox}

\subsection*{Parte 3: Tradúcelo tú}
\begin{codenonum}
mascotas = []
mascotas.append("gato")
mascotas.append("perro")
mascotas.append("pez")
print(len(mascotas))
print(mascotas[1])
\end{codenonum}
Imprime \texttt{3} y luego \texttt{perro}. El elemento 2 de Scratch es el \texttt{[1]} de Python. \py|mascotas = []| hace el trabajo de \chip[lists]{Eliminar todos de mascotas}, porque empieza una lista vacía nuevecita.

\subsection*{Parte 3: Acertijo de dos líneas}
\begin{codenonum}
for palabra in palabras:
    cuenta += 1
\end{codenonum}
El bucle da una vuelta por cada elemento, y en cada vuelta le suma 1 a \texttt{cuenta}. ¡Así es exactamente como cuenta \py|len|!
